%%%%%%%%%%%%%%%%%%%%%%%%%%%%%%
%  Template LaTeX
%  classe: classe_PFC_IFBA_VDC.cls
%  compilador: pdfLaTeX ou XeLaTeX 
%
%  PFC - Projeto de Final de Curso Engenharia Elétrica (COEEL) - IFBA campus Vitória da Conquista e Modelo de Relatório
%
%  Este é um template (modelo) LaTeX que NÃO possui a intenção de seguir a norma ABNT-NBR14724 - Trabalhos Acadêmicos, porém estaremos seguindo as definições da norma ABNT-NBR6023 - Referências
%
%  Template criado para atender às demandas de Projeto Final de Curso (PFC) e Relatórios Técnicos de Engenharia Elétrica do IFBA campus Vitória da Conquista
%
%  versão 2023.08.07
%  Criado por: prof. Elvio Prado da Silva
%  elvio@ieee.org
%  https://github.com/prof-elvio
%
%  distribuição e modificação autorizada desde que mantenha os créditos das pessoas que trabalharam neste template
%%%%%%%%%%%%%%%%%%%%%%%%%%%%%%
\makeatletter % permitir que o símbolo @ seja usado em comandos de macro 
%%%%%%%%%%%%%%%%%%%%%%%%%%%%%%
%%%%%%%%%%%%%%%%%%%%%%%%%%%%%%
%%%%%% DOCUMENTCLASS %%%%%%%%%
%
% [pfc] caso seja PFC - Projeto Final de Curso
% [relatorio] caso seja RELATÓRIO de disciplina
% [preprojeto] caso seja Pré-Projeto de PFC
%
%
\documentclass[pfc]{./classe/cls-PFC-COEEL}
% \documentclass[relatorio]{./classe/cls-PFC-COEEL}
% \documentclass[preprojeto]{./classe/cls-PFC-COEEL}
%
%
%%%%%%%%%%%%%%%%%%%%%%%%%%%%%%
% ABNT-NBR6023 norma para citações bibliográficas
\usepackage[alf]{abntex2cite} 
%
% adicione aqui seus pacotes adicionais
% \usepackage{}
%
%%%%%%%%%%%%%%%%%%%%%%%%%%%%
% Documento que cria a Capa e a folha de Rosto
\input{b_Capa}
%%%%%%%%%%%%%%%%%%%%%%%%%%%%
%%%% CONFIGURA O ÍNDICE REMISSIVO
% configuração das letras em ./classe/estiloindice.ist
\makeindex[columns=3, title=ÍNDICE REMISSIVO, 
           options= -s ./classe/estiloindice.ist] 
%
%%%%%%%%%%%%%%%%%%%%%%%%%%%%
%%%% INICIA O PACOTE PARA GLOSSÁRIO
\loadglsentries{GLOSSARIO.tex}
% se é PFC 
\ifpfc 
    \makeglossaries
\fi 
% se é Pré-Projeto de PFC
\ifpreprojeto 
    \makeglossaries
\fi 
%
%
%%%%%%%%%%%%%%%%%%%%%%%%%%%% 
%%%%%%%%%%%%%%%%%%%%%%%%%%%%
%%%%% BEGIN DOCUMENT %%%%%%%
%%%%%%%%%%%%%%%%%%%%%%%%%%%%
%
%
\begin{document}

%%% 
\input{./classe/inicio}
%%%

% comentar esta linha para seu trabalho
\include{Template} % utilizando o Template
%

%%%%%%%%%%%%%%%%%%%%%%%%%%%
\ifpfc % se é PFC
   \include{./PFC_Especificos/g_Introd_Objetivos} 
   \include{./PFC_Especificos/h_Referencial_Teorico}
   \include{./PFC_Especificos/i_Metodologia}
\fi 
\ifpreprojeto % se é Pré-Projeto
   \include{./PFC_Especificos/g_Introd_Objetivos} 
\fi 
%
%%%%%%%%%%%%%%%%%%%%%%%%%%%%
%%%%%%  AQUI VÃO OS SEUS CAPÍTULOS 
%%%%%%%%%%%%%%%%%%%%%%%%%%%%
% adicione aqui seus capítulos adicionais, adicione quantos forem necessários
% \include{}
% \include{}
\include{Seu_Trabalho}
% \include{Outro_Capitulo}
% \include{Mais_um_Capitulo}
% \include{Mais_outro_Capitulo}
% \include{}
% \include{}
%
%%% não altere este documento 
\input{./classe/fim}

%%%%%%%%%%%%%%%%%%%%%%%%%%%%
%
\end{document}
%
%%%%%%%%%%%%%%%%%%%%%%%%%%%%
%%%%%%%%%%%  FIM  %%%%%%%%%%%%%%%%%%%%%%%%%%%%
%%%%%%%%%%%%%%%%%%%%%%%%%%%%